\documentclass{article}

% Official template target:
% https://www.overleaf.com/latex/templates/neurips-2024/tpsbbrdqcmsh
% Note: as of 2026-05-05, this URL serves the NeurIPS 2026 template and uses
% neurips_2026.sty. Upload this folder into that Overleaf template, or copy the
% official neurips_2026.sty from the template into this directory.
\usepackage{neurips_2026}

\usepackage[utf8]{inputenc}
\usepackage[T1]{fontenc}
\usepackage{hyperref}
\usepackage{url}
\usepackage{booktabs}
\usepackage{amsfonts}
\usepackage{amsmath}
\usepackage{nicefrac}
\usepackage{microtype}
\usepackage{xcolor}
\usepackage{graphicx}
\usepackage{array}

\title{Calibration, Not Prompting, Determines Local LLM Utility on VCBench}

\author{%
  Anonymous Author(s)\\
  Affiliation\\
  Address\\
  \texttt{email}\\
}

\begin{document}

\maketitle

\begin{abstract}
Venture-capital screening is an imbalanced rare-event prediction problem: only a
small fraction of founders produce extreme outcomes, while false positives are
costly because they waste analyst attention and capital. We study this setting
using VCBench, a benchmark for predicting founder success from anonymized
pre-founding profiles. Starting from a prompt-engineering baseline over
Qwen3-32B, we extend the evaluation to local Ollama inference with no-thinking
Qwen models, including \texttt{qwen3:32b} and the quantized
\texttt{hf.co/Qwen/Qwen3-30B-A3B-GGUF:Q4\_K\_M}. Our main finding is that prompt
engineering mostly changes calibration rather than reliably improving
discrimination. On a 120-profile public validation sample, local
\texttt{qwen3:32b} Few-Shot reaches $F_{0.5}=0.1932$ by predicting SUCCESS for
40.8\% of profiles, while local \texttt{qwen3:32b} Vanilla reaches
$F_{0.5}=0.2174$ with a lower but still high 21.7\% predicted-positive rate. In
contrast, the optimized Qwen3-30B-A3B GGUF Vanilla run reaches
$F_{0.5}=0.2632$ with precision 0.5000, but this comes from only two positive
predictions and one true positive. On the 4,500-profile private/blind file, the
same GGUF setup predicts SUCCESS for 333 founders (7.4\%), close to the
benchmark base rate, but labels are unavailable and this is therefore a
submission-style prediction file rather than a scored result. These results
suggest that local open-weight LLMs can be made more conservative, but
small-sample $F_{0.5}$ gains must be interpreted alongside predicted-positive
rate, trivial baselines, and confidence intervals.
\end{abstract}

\section{Introduction}

Early-stage venture-capital prediction is difficult because the signal is
sparse, the labels are noisy, and successful outcomes are rare. A practically
useful screening model should not merely identify many successful founders; it
must do so while keeping false positives low enough to reduce the downstream
review burden. VCBench formalizes this task by asking models to classify
anonymized founder profiles as SUCCESS or FAILURE, with success defined as
raising more than \$500M or achieving an IPO/acquisition above \$500M within
eight years.

The original VCBench framing uses $F_{0.5}$ as the main metric because
precision is more important than recall in capital allocation. However,
$F_{0.5}$ alone can hide important behavior. A model that predicts SUCCESS for
nearly every founder can obtain nonzero recall and a superficially plausible
$F_{0.5}$ while providing little screening utility. Conversely, a model that
predicts very few successes can obtain high precision on a small sample while
missing almost all true positives.

This project began as a prompt-engineering study on Qwen3-32B. The first draft
framed Few-Shot prompting as nearly closing the gap with larger proprietary
models. Subsequent analysis showed that this interpretation was too optimistic:
the apparent gains were driven by overprediction. We therefore revise the
scientific question from ``Can prompting close the frontier-model gap?'' to
``How do prompts and local model variants move the precision-recall and
calibration tradeoff on VCBench?''

This paper reports the latest stage of the work. We evaluate the same prompt
templates under local Ollama inference, disable Qwen thinking mode to reduce
latency and uncontrolled reasoning output, and compare two local Qwen settings.
The core contribution is empirical and methodological: we show that model choice
and inference settings can shift the predicted-positive rate from extreme
overprediction toward the benchmark base rate, but the evidence remains
small-sample and should be judged through calibration diagnostics rather than
$F_{0.5}$ alone.

\paragraph{Contributions.}
We make five contributions. First, we provide a reproducible local Ollama
workflow for VCBench using the same prompt templates as the hosted API
experiments. Second, we document the effect of no-thinking local Qwen inference
on runtime and output parsing. Third, we report public 120-sample validation
results for local \texttt{qwen3:32b} and
\texttt{hf.co/Qwen/Qwen3-30B-A3B-GGUF:Q4\_K\_M}. Fourth, we report the
prediction distribution of a 4,500-row private/blind run for the GGUF model,
while explicitly avoiding unscored private-performance claims. Fifth, we argue
that predicted-positive rate, trivial baselines, and uncertainty estimates
should be mandatory for rare-event LLM classification benchmarks.

\section{Related work}

LLMs have recently been applied to financial decision-making, trading agents,
and future-event forecasting \citep{benhenda2025finrl,xu2025defi,zeng2025futurex,chandak2026future}.
These applications are vulnerable to look-ahead bias when models use information
that would not have been available at the decision time
\citep{benhenda2026lookahead}. VCBench addresses this concern by anonymizing
founder profiles and restricting inputs to pre-founding information
\citep{chen2025vcbench}.

The present work is closest to VCBench and live forecasting benchmarks such as
YCBench \citep{benhenda2026ycbench}. Unlike benchmarks focused primarily on
frontier-model accuracy, this study emphasizes evaluation pathology under class
imbalance. The central issue is not whether an LLM can produce a label, but
whether the label distribution is calibrated enough to provide screening value.

\section{Data and task}

We use the public VCBench split available in this repository:
\texttt{vcbench\_final\_public.csv}. It contains 4,500 anonymized founder
profiles with a success rate of about 9\%. Each row includes
\texttt{anonymised\_prose}, structured background columns, industry, and the
binary label \texttt{success}.

For public validation, we use the existing script protocol: a stratified 80/20
split with seed 42, followed by evaluation on a stratified 120-row sample from
the validation pool. The 120-row sample contains 11 positives and 109 negatives.
Metrics include $F_{0.5}$, precision, recall, TP, FP, FN, TN, and
predicted-success count.

We also run predictions on \texttt{vcbench\_final\_private.csv}. The local copy
used here contains 4,500 rows but no \texttt{success} labels. Therefore, the
private/blind run is a prediction-generation exercise and not a scored
evaluation.

\section{Models and inference}

\subsection{Hosted API baseline}

The original experiment used Qwen3-32B through a Groq OpenAI-compatible API.
That run evaluated four prompt strategies: Vanilla, Chain-of-Thought (CoT),
Few-Shot, and Hybrid. The earlier draft reported that Few-Shot at temperature
0.0 achieved held-out $F_{0.5}=0.1028$ with high recall. Later analysis showed
this result was caused by very high predicted-positive rate and did not clearly
outperform trivial high-positive baselines.

\subsection{Local Ollama inference}

We added local inference scripts, \texttt{run\_experiments\_ollama.py} and
\texttt{test\_real\_data\_ollama.py}. The public local runner imports
\texttt{PROMPTS} directly from \texttt{run\_experiments.py}. The private local
runner imports \texttt{BEST\_PROMPT} directly from \texttt{test\_real\_data.py}.
Thus, the local scripts use the same prompt content as the API scripts; only the
inference backend changes.

Local inference uses Ollama's native chat endpoint at
\texttt{http://127.0.0.1:11434/api/chat}. We use \texttt{127.0.0.1} rather than
\texttt{localhost} because it was more reliable in the Windows environment used
for the experiments.

\subsection{No-thinking mode}

Qwen reasoning models can spend substantial time producing thinking traces.
During local experiments, this made CoT-style prompts especially slow and made
some OpenAI-compatible responses return empty assistant content while reasoning
appeared in a separate field. We therefore use Ollama native chat and disable
thinking with both the \texttt{/no\_think} control token and the JSON option
\texttt{\{"think": false\}}. The final label parser extracts the last occurrence
of SUCCESS or FAILURE from the model output.

\section{Prompt strategies}

The four public prompt strategies are:
\begin{itemize}
    \item \textbf{Vanilla}: direct prediction with the success definition, 9\%
    base-rate warning, and instruction to be conservative.
    \item \textbf{CoT}: explicit positive-signal, risk-factor, and base-rate
    reasoning steps.
    \item \textbf{Few-Shot}: two examples, one failure and one success.
    \item \textbf{Hybrid}: few-shot examples plus structured reasoning.
\end{itemize}
The local full-private script currently uses the same Few-Shot-style
\texttt{BEST\_PROMPT} as \texttt{test\_real\_data.py}.

\section{Results}

\subsection{Public validation: local Qwen3-32B}

The first local Ollama run used \texttt{qwen3:32b}. Full CoT evaluation was
stopped because it was too slow locally; one CoT row took about 75 seconds.
Completed public 120-sample results are shown in Table~\ref{tab:qwen32}.

\begin{table}[t]
\caption{Public 120-sample validation results for local \texttt{qwen3:32b}. The Vanilla prompt reduced false positives relative to Few-Shot and achieved higher $F_{0.5}$.}
\label{tab:qwen32}
\centering
\begin{tabular}{llrrrrrr}
\toprule
Model & Prompt & $T$ & $F_{0.5}$ & Precision & Recall & TP & FP \\
\midrule
\texttt{qwen3:32b} & Few-Shot & 0.0 & 0.1932 & 0.1633 & 0.7273 & 8 & 41 \\
\texttt{qwen3:32b} & Vanilla  & 0.0 & 0.2174 & 0.1923 & 0.4545 & 5 & 21 \\
\bottomrule
\end{tabular}
\end{table}

The Vanilla prompt outperformed Few-Shot in $F_{0.5}$ because it reduced false
positives substantially, even though it found fewer true positives.

\subsection{Public validation: local Qwen3-30B-A3B GGUF}

The next local run used
\texttt{hf.co/Qwen/Qwen3-30B-A3B-GGUF:Q4\_K\_M}. The completed public
validation row is shown in Table~\ref{tab:qwen30}.

\begin{table}[t]
\caption{Public 120-sample validation result for local Qwen3-30B-A3B GGUF. The high precision is based on only two positive predictions, so it should be interpreted as a calibration shift rather than a robust performance claim.}
\label{tab:qwen30}
\centering
\begin{tabular}{llrrrrrrrr}
\toprule
Model & Prompt & $T$ & $F_{0.5}$ & Precision & Recall & TP & FP & FN & Pred+ \\
\midrule
Qwen3-30B-A3B GGUF & Vanilla & 0.0 & 0.2632 & 0.5000 & 0.0909 & 1 & 1 & 10 & 1.7\% \\
\bottomrule
\end{tabular}
\end{table}

This is the highest $F_{0.5}$ observed in the latest local public sample, but
it should be interpreted carefully. The result depends on only two positive
predictions and one true positive. It is evidence of a conservative calibration
shift, not yet evidence of robust superiority.

\subsection{Private/blind prediction distribution}

The full private/blind file contains 4,500 profiles and no labels. The local
GGUF run produced the prediction distribution in Table~\ref{tab:private_dist}.

\begin{table}[t]
\caption{Prediction distribution on the 4,500-row private/blind file. Labels are unavailable, so this is not a scored result.}
\label{tab:private_dist}
\centering
\begin{tabular}{lrr}
\toprule
Prediction & Count & Rate \\
\midrule
SUCCESS & 333  & 7.4\% \\
FAILURE & 4167 & 92.6\% \\
\bottomrule
\end{tabular}
\end{table}

This distribution is close to the public benchmark base rate of about 9\%.
However, without labels, we cannot compute $F_{0.5}$, precision, recall, or
compare against official private baselines.

\section{Discussion}

The latest results support three conclusions. First, prompt engineering alone is
not the main story. The same family of prompts can produce very different
predicted-positive rates depending on model backend, quantization, and inference
settings. This makes calibration diagnostics essential.

Second, no-thinking local inference is practically important. It reduces
uncontrolled reasoning output and makes local evaluation more feasible. It also
prevents a mismatch where the model spends the token budget on reasoning while
returning no usable classification content.

Third, the GGUF model appears more conservative than the earlier local
\texttt{qwen3:32b} setting. This is promising for VC screening because extreme
overprediction is harmful. But conservatism can also become underprediction: the
public validation row recalls only 1 of 11 successes. A useful system must
balance this tradeoff across larger samples.

\section{Limitations}

The public validation results are based on a 120-row sample. The strongest GGUF
public result has only one true positive and one false positive. CoT and Hybrid
were not fully rerun under the latest local GGUF setup because local reasoning
prompts were too slow for practical iteration. The private/blind predictions are
unscored because labels are unavailable. The study has not yet evaluated
multiple seeds, confidence intervals, or full public-split bootstrap estimates.
Finally, the structured-feature ML baselines discussed in the earlier draft need
to be fully reproduced and included in the repository before they can support a
publication claim.

\section{Broader impact}

Automated VC screening can reinforce existing inequities if models over-weight
prestige signals such as elite universities, prior big-tech employment, or
geographic access to startup ecosystems. Even anonymized profiles can retain
strong socioeconomic proxies. The present work should therefore be interpreted
as benchmark research, not as a deployable investment system. Any deployment
would require fairness audits, human oversight, calibration monitoring, and a
clear policy that model outputs are decision support rather than automated
funding decisions.

\section{Conclusion}

The latest experiments shift the paper's conclusion from ``prompt engineering
can make open-weight LLMs competitive on VCBench'' to a more defensible claim:
calibration determines whether local LLMs are useful for rare-event VC
screening. Local no-thinking Qwen inference can move predictions from severe
overprediction toward much more conservative behavior, and the current GGUF
private/blind run predicts SUCCESS at a rate close to the public base rate.
However, the evidence is not yet strong enough for a final performance claim.

For a NeurIPS-quality submission, the next required step is a larger and more
statistically grounded evaluation: full public split or repeated stratified
splits, trivial baselines in every table, confidence intervals, and a fully
reproducible structured-ML baseline. With those additions, the work could be a
credible negative-result and evaluation-methodology paper for LLMs on
imbalanced, high-stakes prediction tasks.

\begin{ack}
This draft was prepared from local VCBench experiments and documentation in the
associated repository. No external funding or competing interests are declared
in this draft version.
\end{ack}

\bibliographystyle{plainnat}
\bibliography{references}

\appendix

\section{Reproducibility details}

Code for local inference is provided in \texttt{run\_experiments\_ollama.py} and
\texttt{test\_real\_data\_ollama.py}. Public validation uses
\texttt{vcbench\_final\_public.csv}. The private/blind prediction run uses
\texttt{vcbench\_final\_private.csv}. The current local model is
\texttt{hf.co/Qwen/Qwen3-30B-A3B-GGUF:Q4\_K\_M}. The inference backend is the
Ollama native chat endpoint, with no-thinking controls \texttt{/no\_think} and
\texttt{think: false}. The public validation sample size is 120, with 11
positives, using seed 42.

\newpage
\section*{NeurIPS Paper Checklist}

\begin{enumerate}
\item \textbf{Claims}

Question: Do the main claims made in the abstract and introduction accurately
reflect the paper's contributions and scope?

Answer: \answerYes{}.

Justification: The abstract and introduction state that the work is a
small-sample calibration study and explicitly avoid claiming scored private-set
performance.

\item \textbf{Limitations}

Question: Does the paper discuss the limitations of the work?

Answer: \answerYes{}.

Justification: Section 8 discusses sample size, unscored private predictions,
missing confidence intervals, incomplete CoT/Hybrid reruns, and incomplete
structured-feature baselines.

\item \textbf{Theory assumptions and proofs}

Question: For each theoretical result, does the paper provide the full set of
assumptions and a complete proof?

Answer: \answerNA{}.

Justification: The paper is empirical and does not present new theoretical
results or proofs.

\item \textbf{Experimental reproducibility}

Question: Does the paper fully disclose the information needed to reproduce the
main experimental results?

Answer: \answerYes{}.

Justification: The paper reports datasets, scripts, model identifiers, sample
size, seed, prompts, Ollama endpoint, and no-thinking controls. The repository
contains the corresponding code and result CSVs.

\item \textbf{Open access to data and code}

Question: Does the paper provide open access to the data and code, with
sufficient instructions to reproduce the results?

Answer: \answerYes{}.

Justification: The public data and scripts are in the repository. The
private/blind file is used only for prediction generation and is not presented as
a scored benchmark.

\item \textbf{Experimental setting details}

Question: Does the paper specify all training and evaluation details needed to
understand the results?

Answer: \answerYes{}.

Justification: Sections 3--6 describe the split, sample size, model identifiers,
prompt strategies, decoding setup, output parser, and metrics.

\item \textbf{Statistical significance}

Question: Does the paper report uncertainty estimates or statistical
significance where appropriate?

Answer: \answerNo{}.

Justification: The paper explicitly identifies missing confidence intervals and
multi-seed evaluation as a limitation and next required step.

\item \textbf{Datasets}

Question: Does the paper describe the dataset characteristics and any relevant
privacy or licensing constraints?

Answer: \answerYes{}.

Justification: Section 3 describes the public and private/blind VCBench files,
the 9\% class imbalance, the missing private labels, and the anonymized founder
profile setting.

\item \textbf{Compute resources}

Question: Does the paper report the compute resources used?

Answer: \answerNo{}.

Justification: The current draft reports local Ollama inference and qualitative
runtime observations, but it does not yet include hardware specifications. This
should be added before submission.

\item \textbf{Broader impacts}

Question: Does the paper discuss both potential positive and negative societal
impacts?

Answer: \answerYes{}.

Justification: Section 9 discusses deployment risks, fairness concerns, prestige
proxies, and the need for human oversight in VC screening.

\item \textbf{Safeguards}

Question: Does the paper describe safeguards for responsible use where relevant?

Answer: \answerYes{}.

Justification: The paper frames the system as benchmark research, not a
deployment-ready screening system, and recommends fairness audits, calibration
monitoring, and human oversight.

\item \textbf{Licenses for existing assets}

Question: Are licenses for existing assets such as code, data, and models
identified where possible?

Answer: \answerNo{}.

Justification: The draft names the model and benchmark files, but model and
dataset license information should be verified and added before submission.
\end{enumerate}


\end{document}
